% =====================================================================
%  DISSERTATION — complete, single file
%
%  Compile from the repository root (figures are read from
%  results/figures/*.pdf):
%
%      pdflatex dissertation
%      bibtex   dissertation
%      pdflatex dissertation
%      pdflatex dissertation
%
%  If IEEEtran.bst is unavailable, change \bibliographystyle to plain.
% =====================================================================
\documentclass[12pt,a4paper,oneside]{report}

\usepackage[a4paper,left=3.5cm,right=2.5cm,top=2.5cm,bottom=2.5cm]{geometry}
\usepackage{setspace}
\onehalfspacing

\usepackage[T1]{fontenc}
\usepackage[utf8]{inputenc}
\IfFileExists{lmodern.sty}{\usepackage{lmodern}}{}
\usepackage{textcomp}

\usepackage{graphicx}
\usepackage{booktabs}
\usepackage{array}
\usepackage{longtable}
\usepackage{amsmath,amssymb}
\usepackage{enumitem}
\usepackage{xcolor}
\usepackage{listings}
\usepackage{caption}
\usepackage{url}
\usepackage[hidelinks]{hyperref}

\definecolor{codebg}{RGB}{248,248,250}
\definecolor{codekw}{RGB}{18,42,86}
\definecolor{codestr}{RGB}{22,110,72}
\definecolor{codecom}{RGB}{120,126,134}
\definecolor{good}{RGB}{22,110,72}
\definecolor{bad}{RGB}{158,42,42}

\lstdefinestyle{py}{
  language=Python, basicstyle=\ttfamily\footnotesize,
  keywordstyle=\color{codekw}\bfseries, stringstyle=\color{codestr},
  commentstyle=\color{codecom}\itshape, backgroundcolor=\color{codebg},
  frame=leftline, framerule=1.2pt, rulecolor=\color{codekw!40},
  xleftmargin=10pt, framexleftmargin=8pt, breaklines=true,
  showstringspaces=false, aboveskip=8pt, belowskip=8pt, captionpos=b,
}
\lstdefinestyle{plainlist}{
  language={}, basicstyle=\ttfamily\footnotesize,
  backgroundcolor=\color{codebg}, frame=leftline, framerule=1.2pt,
  rulecolor=\color{codekw!40}, xleftmargin=10pt, framexleftmargin=8pt,
  breaklines=true, showstringspaces=false, aboveskip=8pt, belowskip=8pt,
}
\lstset{style=py}

\captionsetup{font=small, labelfont=bf, skip=6pt}

\usepackage{titlesec}
\titleformat{\chapter}[display]
  {\normalfont\huge\bfseries}{\chaptertitlename\ \thechapter}{18pt}{\Large}
\titlespacing*{\chapter}{0pt}{-20pt}{28pt}

\newcommand{\up}[1]{\textcolor{good}{\textbf{#1}}}
\newcommand{\down}[1]{\textcolor{bad}{\textbf{#1}}}

\begin{document}

% =====================================================================
\begin{titlepage}
\centering
\vspace*{1.2cm}
{\LARGE\bfseries Spatial Mental Modeling from Limited Views\par}
\vspace{0.8cm}
{\Large Reproducing ViperGPT without Codex:\\[3pt]
a Prompt-Specification Gap and a Depth-Grounded API\par}
\vspace{2.0cm}
{\large A Dissertation submitted in partial fulfilment\\
of the requirements for the degree of\par}
\vspace{0.5cm}
{\large\bfseries Master of Technology\par}
{\large in\par}
{\large\bfseries Computer Science and Engineering\\
(Information Security and Privacy)\par}
\vspace{1.8cm}
{\large\itshape submitted by\par}
\vspace{0.25cm}
{\Large\bfseries MD Humayun\par}
\vspace{1.4cm}
{\large\itshape under the supervision of\par}
\vspace{0.25cm}
{\Large\bfseries Dr.\ M.\ A.\ Zaveri\par}
\vfill
{\large Department of Computer Science and Engineering\\
Sardar Vallabhbhai National Institute of Technology\\
Surat, Gujarat, India\par}
\vspace{0.6cm}
{\large \the\year\par}
\end{titlepage}

% =====================================================================
\pagenumbering{roman}
\setcounter{page}{2}

\chapter*{Abstract}
\addcontentsline{toc}{chapter}{Abstract}

ViperGPT answers a visual query by prompting a code-generating language model to emit
a Python program over pretrained vision modules, and then executing that program. Its
published generator, OpenAI Codex, has since been withdrawn, so any present-day
reproduction must substitute a different model. This dissertation reports what that
substitution reveals.

Rebuilding the pipeline with an open-weights generator produces an accuracy of
$0.00$ on RefCOCO referring-expression comprehension. The failure is not a crash:
$332$ of $500$ programs execute to completion and return a natural-language string
where a bounding box is required. The released API specification never states the
return type a grounding program must produce, and five of its seven worked examples
declare \texttt{str}. Stating the contract explicitly raises accuracy to $39.20$.
The effect is independent of model scale --- it appears at 1.5B, 7B and 32B
parameters and is repaired at each --- and scaling the generator twentyfold buys
approximately six accuracy points where stating the contract buys thirty-three to
forty-one.

A second finding follows. Across those three model sizes, non-spatial accuracy rises
from $33.65$ to $48.82$ while spatial accuracy remains flat at $34.95$, $34.26$ and
$34.95$, on identical images with identical perception modules. A larger generator
writes better programs for everything except spatial relations. The bottleneck is the
interface: the API exposes a single scalar for depth, from which inter-object
ordering cannot be expressed, and the prompt directs the generator away from geometry.

The dissertation therefore extends the API with depth-ordering primitives and a
corrected depth backend. Offered without qualification these improve depth-word
queries from $4.26$ to $19.15$ but degrade non-spatial accuracy by ten points: given
a new primitive, the generator applies it indiscriminately. Adding three sentences
that delimit when the primitives apply improves every subset of both datasets,
raising depth queries to $21.28$ and overall RefCOCO accuracy to $42.00$
($42.07 \pm 0.61$ over three seeds against $36.07 \pm 1.75$ for the contract alone),
while reducing the seed-to-seed variance of spatial queries by a factor of $4.6$.

Two results large enough to be attractive were rejected after control experiments.
Capping the detector's candidate set raises accuracy by fifteen points, but a
condition containing no depth primitives gains almost as much, so the improvement
belongs to the candidate set rather than to the method. Separately, a depth-sign
correction that a synthetic validation had rejected was reinstated after testing
against ground-truth annotations, where the uncorrected signal scored below chance.

\vspace{1em}
\noindent\textbf{Keywords:} visual reasoning, program synthesis, referring expression
comprehension, reproducibility, spatial grounding, vision--language models.

\tableofcontents
\listoffigures
\listoftables

\cleardoublepage
\pagenumbering{arabic}
\setcounter{page}{1}

% =====================================================================
\chapter{Introduction}
\label{ch:introduction}

\section{Background}

A question asked about an image is rarely atomic. \emph{``The man on the right''}
requires locating every man in the scene and then ordering them; \emph{``the cup
behind the bottle''} requires locating two object classes and then comparing their
distance from the camera. Monolithic vision--language models answer such queries in a
single forward pass, which makes the intermediate reasoning inaccessible: when the
answer is wrong, there is no step to point at.

An alternative treats the query as a \emph{program to be written}. A large language
model is shown a specification of an API over pretrained vision modules and asked to
emit code that composes those modules; the code is then executed against the image.
Neural Module Networks introduced composition over learned modules
\cite{andreas2016nmn}, and capable code-generation models made it possible to
synthesise the composition directly in a general-purpose language. VisProg
\cite{gupta2023visprog} and ViperGPT \cite{suris2023vipergpt} are the prominent
instances.

The appeal is threefold. The system is \emph{training-free}: no gradient step is
taken, and new capabilities are added by extending the API rather than by retraining.
It is \emph{compositional}: a query never seen during training can be answered if its
parts can be. And it is \emph{interpretable} in a strong sense, since the intermediate
reasoning is a Python program a human can read and a machine can execute step by step.

\section{ViperGPT}

ViperGPT formalises this as a two-stage pipeline. A code generator $\pi$ maps a query
$q$ to a program $z$, and an execution engine $\phi$ runs that program against an
image $x$:
\[
    z = \pi(q), \qquad r = \phi(x, z).
\]
The generator is prompted with an API specification describing an \texttt{ImagePatch}
class and functions over it --- \texttt{find} for open-vocabulary detection,
\texttt{verify\_property} for attributes, \texttt{simple\_query} for visual question
answering, \texttt{compute\_depth} for relative depth, and others. The generator never
sees the image; it sees only the query and the API. Execution binds the API to
pretrained models: GLIP \cite{li2022glip} for detection, X-VLM \cite{zeng2022xvlm}
for attributes, BLIP-2 \cite{li2023blip2} for question answering, MiDaS
\cite{ranftl2020midas} for depth.

On referring-expression comprehension the paper reports $72.0$ on RefCOCO and $67.0$
on RefCOCO+ (testA), against $55.0$ and $52.2$ for GLIP used directly and $58.6$ and
$60.5$ for ReCLIP \cite{subramanian2022reclip}.

\section{Problem statement}

The generator used in that work, \texttt{code-davinci-002}, has been withdrawn. Any
attempt to reproduce or build on ViperGPT today must substitute a different generator,
and no faithful option exists. This is not a peculiarity of one paper: a published
system whose central component is reachable only through a proprietary API is
reproducible only while that API is maintained.

Substitution raises a question the original work had no reason to ask. \emph{Which
parts of the reported behaviour belong to the method, and which belonged to Codex?}
An API specification written for one model is an interface contract; if part of that
contract was carried implicitly by the model's own inference rather than by the text,
the omission is invisible until the model changes.

A second question follows from the API itself. Of the operations ViperGPT exposes,
exactly one concerns three-dimensional structure: \texttt{compute\_depth()}, returning
a single scalar. From one scalar per object it is possible to express \emph{``A is
behind B''}. It is not possible to express \emph{``the third-closest person''},
\emph{``between A and B in depth''}, or physical size. Meanwhile the prompt directs
the generator to resolve \emph{left}, \emph{right}, \emph{up} and \emph{down} using
base Python on patch coordinates. Whether this suffices for the spatial language
referring expressions actually contain is an empirical question that an aggregate
number per dataset does not answer.

\section{Research questions}

\begin{description}[leftmargin=2.2em, style=nextline, itemsep=4pt]
\item[RQ1] When Codex is replaced by an open-weights code generator, what happens on
referring-expression comprehension, and how much of any shortfall is attributable to
the released prompt rather than to the substituted model?
\item[RQ2] Is that effect a property of the prompt or of the generator's capacity ---
does it change with model scale?
\item[RQ3] Does accuracy differ systematically by the \emph{kind} of spatial language
a query contains, and is any difference explained by the expressiveness of the API?
\item[RQ4] Does extending the API with depth-ordering primitives improve accuracy on
the queries that need them, without degrading those that do not?
\end{description}

\section{Approach}

The work proceeds in three milestones, each stated as a hypothesis and tested.

\textbf{Milestone 1} reproduces the generation stage alone. Programs are synthesised
for 500 RefCOCO and 500 RefCOCO+ queries and analysed \emph{statically} --- parsed,
not executed --- so the generator can be studied before the perception stack exists.
The initial hypothesis, that spatial queries would collapse into a single opaque
\texttt{simple\_query} call, was rejected; the analysis instead revealed a return-type
failure affecting $98.8\%$ of programs.

\textbf{Milestone 2} builds the execution stack and measures end-to-end accuracy.
This required compiling the GLIP fork for the Hopper architecture, which the published
environment cannot target, and substituting CLIP \cite{radford2021clip} for X-VLM.
A hypothesis that RefCOCO+ would provide a clean control for spatial language was
rejected: the gap widened rather than narrowed, and a manual audit established that
the two datasets contain different \emph{kinds} of spatial language.

\textbf{Milestone 3} extends the API with depth-grounded primitives. A first form,
offering the primitives without specifying when they apply, improved the target
queries and degraded the others; the final form adds a routing rule and improves every
subset measured.

\section{Contributions}

\begin{enumerate}[leftmargin=2.0em, itemsep=5pt]
\item \textbf{A working reproduction with an open-weights generator}, on hardware the
published environment cannot run on. The GLIP fork is rebuilt for compute capability
9.0 and verified independently of the full pipeline.

\item \textbf{Identification and quantification of a specification gap.} The released
prompt does not state the required return type, and five of its seven worked examples
declare \texttt{str}. With an open-weights generator, $98.8\%$ of programs return a
string where a box is required; executed, this yields $0.00$. Stating the contract
raises it to $39.20$.

\item \textbf{Evidence that this is a property of the specification, not of capacity.}
The failure occurs at 1.5B, 7B and 32B ($95.8\%$, $98.8\%$, $87.6\%$ task-invalid) and
is repaired at every scale. Scaling twentyfold buys $\approx 6$ points; the contract
buys $33$--$41$.

\item \textbf{Evidence that spatial accuracy does not respond to scale.} Non-spatial
accuracy rises $33.65 \to 48.82$ across the three models while spatial accuracy stays
at $34.95 / 34.26 / 34.95$ on identical images and modules.

\item \textbf{A depth-grounded API extension} --- \texttt{depth\_order},
\texttt{is\_behind}, \texttt{is\_in\_front\_of}, \texttt{is\_between\_3d},
\texttt{distance\_3d} --- with a corrected depth backend. The generator adopts it,
invoking \texttt{depth\_order} $119$ times across $500$ programs where it invokes the
original \texttt{compute\_depth} $29$ times.

\item \textbf{The finding that providing a capability is insufficient without
specifying when it applies.} The primitives alone move depth queries $4.26 \to 19.15$
and cost ten points of non-spatial accuracy; three sentences delimiting their scope
improve every subset, reaching $21.28$ on depth queries and $42.00$ overall.

\item \textbf{Two negative results established by control rather than omitted}: a
15-point apparent gain from capping the candidate set, shown to arise from the
candidate set rather than from depth reasoning; and a depth-sign error that a
synthetic validation had endorsed, identified by testing against ground truth.
\end{enumerate}

\section{Scope}

Evaluation is restricted to referring-expression comprehension on RefCOCO and
RefCOCO+ --- Table~1 of the original paper. GQA and NExT-QA are excluded for storage
and compute; OK-VQA depends on GPT-3, likewise withdrawn. No cross-task claim is made.

The reproduction is not faithful and is not presented as such: the final configuration
reaches $42.00$ against a reported $72.0$, and the responsible deviations are
enumerated in Chapter~\ref{ch:deviations}. What is reproduced is the \emph{method};
what is measured is what happens to it when its generator is replaced.

Depth is claimed only as \emph{ordering}. The datasets provide no camera intrinsics,
so metric distances would rest on an assumed field of view; the ordering primitives
require only monotonicity in true distance, which relative depth provides once its
sign is correct.

% =====================================================================
\chapter{Literature Review}
\label{ch:literature}

\section{Compositional visual reasoning}

The observation that a visual question decomposes into simpler operations predates
large language models. Neural Module Networks \cite{andreas2016nmn} parse a question
into a layout of learned modules and assemble a network per instance. The modules are
trained jointly, so the approach inherits both the strengths of end-to-end learning
and its costs: a new capability requires retraining, and the layout parser must itself
be learned.

Later work refined the parsing rather than the principle. Stack-NMN
\cite{hu2018stacknmn} makes the layout differentiable; Neural-Symbolic VQA
\cite{yi2018nsvqa} separates perception from a symbolic executor over a structured
scene representation. The recurring difficulty is that the program space must be fixed
in advance, because it is the space over which the parser is trained.

\section{Program synthesis with language models}

Code-generating models change that constraint. If a model can write Python from a
specification, the program space is the language itself, and adding a capability
becomes a matter of documenting a function rather than retraining a parser.

VisProg \cite{gupta2023visprog} generates programs in a small domain-specific language
over a fixed module set. ViperGPT \cite{suris2023vipergpt} generates ordinary Python,
supplies an API specification as the prompt, and executes the result with
\texttt{exec}. CodeVQA \cite{subramanian2023codevqa} applies the idea to question
answering with a smaller module set; Chameleon \cite{lu2023chameleon} generalises to
tool sequences beyond vision.

These systems share three properties: they are training-free, their intermediate
representation is inspectable, and their behaviour is determined by two artefacts that
receive little attention in evaluation --- the API surface and the prompt describing
it. This dissertation studies the second pair.

\subsection{The generator as an undocumented dependency}

ViperGPT uses Codex; VisProg uses GPT-3; CodeVQA uses Codex. All three models have
been withdrawn. The published prompts remain, but the model that interpreted them does
not, and there is no way to establish after the fact which parts of a system's
behaviour were specified by the prompt and which were supplied by the model.

Related concerns appear in the reproducibility literature. Kapoor and Narayanan
\cite{kapoor2023leakage} document how evaluation flaws propagate through published
results; Bell and Kampman \cite{bell2021perspectives} survey reproducibility in AI
research more broadly. The failure mode considered here --- a system whose central
component becomes permanently unavailable --- is a consequence of building on closed
APIs, and this work is a case study in what such a dependency conceals.

\section{Referring expression comprehension}

Referring expression comprehension asks a model to localise the object described by a
phrase. RefCOCO and RefCOCO+ \cite{yu2016refcoco} are the standard benchmarks,
collected on COCO images. They differ in one deliberate respect: RefCOCO permits
location words, RefCOCO+ forbids them, so its expressions describe objects by
appearance rather than position.

Strong baselines follow two routes. Detection-based methods --- MDETR
\cite{kamath2021mdetr}, GLIP \cite{li2022glip}, OWL-ViT \cite{minderer2022owlvit} ---
pretrain for grounded detection directly. Composition-based methods such as ReCLIP
\cite{subramanian2022reclip} decompose the expression and score candidate regions with
CLIP without task-specific training. ViperGPT belongs to a third category: it neither
trains nor scores directly, but writes a program that calls a detector.

\subsection{The spatial component}

An aspect of RefCOCO that bears on this work is that its expressions contain several
distinct kinds of spatial language. \emph{Left}, \emph{right}, \emph{top} and
\emph{bottom} are viewer-relative and resolvable in the image plane. \emph{Closest},
\emph{farthest}, \emph{behind} and \emph{in front of} concern distance from the camera
and are not. Because RefCOCO+ excludes location words but not distance words, the two
datasets differ in the \emph{composition} of their spatial subsets, not merely in
quantity. Aggregate accuracy cannot distinguish these cases;
Chapter~\ref{ch:results} shows they behave differently.

\section{Spatial reasoning in vision--language models}

Recent work reports that vision--language models handle spatial relations poorly.
SpatialVLM \cite{chen2024spatialvlm} attributes this to a lack of 3D supervision and
synthesises spatial question-answering data at scale. Benchmarks such as What's Up
\cite{kamath2023whatsup} and VSR \cite{liu2023vsr} isolate spatial relations and find
large models near chance on some of them.

That literature concerns models trained end to end, where the deficit is a property of
the learned representation. In a program-synthesis system the situation differs:
spatial relations are resolved by code, so the deficit may lie in the model, in the
modules, or in the interface between them. Distinguishing these can be done
experimentally, and Chapter~\ref{ch:results} does so by varying model scale while
holding the interface fixed.

\section{Positioning}

This dissertation is a reproduction study with an extension. It differs from the above
in three respects. It treats the API specification as an object of study rather than
as infrastructure. It separates the contribution of the generator from that of the
interface by varying model scale --- a control available in a training-free system and
rarely reported. And it reports two negative results established by control rather
than omitting them.

% =====================================================================
\chapter{Problem Formulation}
\label{ch:problem}

\section{Formal setting}

Let $q$ be a query, $x$ an image, and $\mathcal{A}$ an API consisting of functions
over image regions together with a textual specification $s(\mathcal{A})$. A generator
$\pi$ maps specification and query to a program, and an engine $\phi$ evaluates it:
\[
    z \;=\; \pi\bigl(s(\mathcal{A}),\, q\bigr),
    \qquad
    r \;=\; \phi\bigl(x,\, z;\, \mathcal{A}\bigr).
\]
For grounding the required output is a region, so a correct program satisfies
$r \in \mathcal{B}$. Accuracy is measured by intersection over union against the
ground-truth box $b^{*}$:
\[
    \mathrm{IoU}(r, b^{*}) = \frac{\lvert r \cap b^{*}\rvert}{\lvert r \cup b^{*}\rvert},
    \qquad \text{correct} \iff \mathrm{IoU}(r, b^{*}) \geq \tau, \quad \tau = 0.5 .
\]

Three properties matter. The generator never observes $x$. Nothing in the type system
enforces $r \in \mathcal{B}$, because $z$ is ordinary Python and $\phi$ is
\texttt{exec}. And $s(\mathcal{A})$ is prose, not a schema, so any constraint it fails
to state is unenforced.

\section{Problem 1: the generator is unavailable}

ViperGPT instantiates $\pi$ as Codex, which has been withdrawn; reproduction requires
$\pi'$. The substitution changes the entity that interprets $s(\mathcal{A})$. If some
constraint $c$ holds of the programs Codex produced, two explanations exist: $c$ is
stated in $s(\mathcal{A})$, or Codex inferred it. These are indistinguishable while the
model is available and distinguished the moment it is replaced.

\section{Problem 2: the specification is incomplete}

The released specification is task-agnostic --- the same text serves question
answering and grounding. It contains seven worked examples, whose declared return
types are distributed as in Table~\ref{tab:returntypes}.

\begin{table}[htbp]\centering
\caption{Return types declared by the worked examples in the released specification
(\texttt{prompts/api.prompt}, 7{,}645 bytes). Line numbers refer to that file.}
\label{tab:returntypes}
\begin{tabular}{@{}llc@{}}
\toprule
\textbf{Example documents} & \textbf{Declared return type} & \textbf{Line} \\
\midrule
\texttt{find}             & \texttt{List[ImagePatch]} & 44 \\
\texttt{exists}           & \texttt{str} & 55 \\
\texttt{verify\_property} & \texttt{str} & 68 \\
\texttt{best\_text\_match}& \texttt{str} & 79 \\
\texttt{simple\_query}    & \texttt{str} & 92 \\
\texttt{compute\_depth}   & \texttt{ImagePatch} & 104 \\
\texttt{llm\_query}       & \texttt{str} & 123 \\
\bottomrule
\end{tabular}
\end{table}

Five of seven return a string. One returns a list. Exactly one returns the type a
grounding program requires, and it sits inside the docstring of
\texttt{compute\_depth}, where a reader seeking the output contract has no reason to
look. No sentence states that \texttt{execute\_command} must return an
\texttt{ImagePatch}.

\section{Problem 3: the interface cannot express depth relations}

Let $d(p)$ denote \texttt{compute\_depth} for a patch $p$. The relation \emph{$a$ is
behind $b$} is expressible as $d(a) > d(b)$. The following are not expressible without
the program implementing the ordering logic itself:

\begin{itemize}[leftmargin=1.6em, itemsep=2pt]
\item \emph{the third-closest person} --- a total order over a set;
\item \emph{the cup between the bottle and the plate in depth} --- an interval test
over three patches;
\item \emph{the larger of the two objects} --- physical size, hence back-projection.
\end{itemize}

The specification compounds this. Among its three guidelines:

\begin{quote}\itshape
Use base Python (comparison, sorting) for basic logical operations,
left/right/up/down, math, etc.
\end{quote}

This is correct for image-plane relations and cannot be correct for depth relations,
because the image plane does not contain the information.

\section{Problem 4: whether an extended interface helps}

If the interface is the constraint, extending it should improve the affected queries.
Two failure modes must be distinguished from success: the generator may ignore the new
functions, as it largely ignores \texttt{compute\_depth}; or it may apply them
indiscriminately, improving the target queries at the expense of others.

\section{Criteria for a claim}

Because this is a reproduction study, the standard of evidence is fixed in advance.

\begin{description}[leftmargin=2.2em, style=nextline, itemsep=4pt]
\item[C-1 Attribution] Any improvement must be shown to arise from the intervention
rather than from a change in the evaluation. Where an intervention could plausibly help
for an unrelated reason, a control condition without it is run.
\item[C-2 Non-degradation] An improvement confined to a target subset at the cost of
others is reported as a trade-off, not as a gain.
\item[C-3 Stability] Where decoding is stochastic, results are reported as mean and
standard deviation over three seeds.
\item[C-4 Provenance] Every reported figure traces to a stored run summary, a
configuration, and a scheduler job identifier.
\item[C-5 Disclosure] Every departure from the published configuration is recorded with
its reason and expected effect, whether or not it is believed to matter.
\end{description}

Criteria C-1 and C-2 are not decorative. Chapter~\ref{ch:rejected} reports two
interventions rejected under C-1, and one that satisfied its target subset but failed
C-2 and was superseded.

% =====================================================================
\chapter{System Architecture and Methodology}
\label{ch:methodology}

\section{Architecture}

Figure~\ref{fig:architecture} shows the reproduced pipeline. Colour marks provenance:
components carried over unchanged, components substituted because the published choice
is unavailable, and components introduced by this work.

\begin{figure}[htbp]\centering
\includegraphics[width=0.80\textwidth]{results/figures/fig14_architecture.pdf}
\caption[Reproduced system architecture]{The reproduced pipeline. A query and an API
specification are given to the program generator; the emitted program is written to
disk, then executed against the image by an engine that binds the API to pretrained
modules. The predicted region is scored by IoU.}
\label{fig:architecture}
\end{figure}

\subsection{The stage boundary}

The published design runs generation and execution in one process. This implementation
writes programs to \texttt{programs.jsonl} between the stages.

The immediate reason is a dependency conflict: the GLIP fork requires an older PyTorch
than the code generator, and the two cannot share an interpreter. Two consequences
follow, both useful. Prompt conditions are compared on \emph{identical} programs, since
the execution stage reads a fixed file. And the generation stage can be studied before
the execution stage exists, which is what makes the static analysis of
Section~\ref{sec:static} possible.

\section{Departures from the published design}

\begin{figure}[htbp]\centering
\includegraphics[width=\textwidth]{results/figures/fig15_paper_vs_ours.pdf}
\caption[Published design against this implementation]{Component-by-component
comparison. Amber marks a substitution forced by unavailability; teal marks a
deliberate change. The unchanged elements are the method, the executor semantics, the
detector, and the task.}
\label{fig:paperVSours}
\end{figure}

Every departure carries an identifier D1--D12 and is documented in
Chapter~\ref{ch:deviations}. The substitutions that most affect results are the
generator (D1) and the attribute model (D9).

\section{Prompt specification}
\label{sec:prompts}

The experimental variable is the API specification. Four variants are defined; each is
a strict superset of the previous one, so any difference in outcome is attributable to
the added text. Model, modules, data and decoding are held fixed.

\begin{table}[htbp]\centering
\caption{The four prompt conditions.}
\label{tab:conditions}
\begin{tabular}{@{}llrp{0.40\textwidth}@{}}
\toprule
& \textbf{File} & \textbf{Bytes} & \textbf{Added} \\
\midrule
C1 & \texttt{api.prompt} & 7{,}645 & --- the released specification \\
C2 & \texttt{api\_grounding.prompt} & 8{,}639 & an explicit output contract, two
grounding examples \\
C3 & \texttt{api\_depth.prompt} & 11{,}485 & five depth primitives with docstrings and
three examples; removal of the instruction directing spatial relations to base Python \\
C4 & \texttt{api\_depth\_cond.prompt} & 12{,}680 & three routing rules delimiting when
the primitives apply, and two counter-examples \\
\bottomrule
\end{tabular}
\end{table}

\subsection{C1 --- the released specification}

The prompt declares an \texttt{ImagePatch} class with attributes \texttt{left},
\texttt{right}, \texttt{upper}, \texttt{lower}, \texttt{horizontal\_center},
\texttt{vertical\_center}, and methods \texttt{find}, \texttt{exists},
\texttt{verify\_property}, \texttt{best\_text\_match}, \texttt{simple\_query},
\texttt{compute\_depth}, \texttt{crop}, \texttt{overlaps\_with}, together with free
functions \texttt{best\_image\_match}, \texttt{distance}, \texttt{bool\_to\_yesno} and
\texttt{coerce\_to\_numeric}. It closes with:

\begin{lstlisting}[style=plainlist]
Consider the following guidelines:
- Use base Python (comparison, sorting) for basic logical operations,
  left/right/up/down, math, etc.
- Use the llm_query function to access external information and answer
  informational questions not concerning the image.
- Return only the function; do not include explanations or example usage.

Query: INSERT_QUERY_HERE
\end{lstlisting}

Two properties are load-bearing. The text nowhere states the required return type, and
five of seven examples declare \texttt{str}. And the first guideline directs spatial
relations to base Python, excluding the depth API by construction.

The single example returning a patch also fixes the semantics of
\texttt{compute\_depth}:

\begin{lstlisting}[caption={The only example in the released prompt that returns an
\texttt{ImagePatch}, quoted verbatim from \texttt{prompts/api.prompt}, lines 100--110.}]
def compute_depth(self):
    """Returns the median depth of the image crop.
    Examples
    --------
    >>> # the bar furthest away
    >>> def execute_command(image)->ImagePatch:
    >>>     image_patch = ImagePatch(image)
    >>>     bar_patches = image_patch.find("bar")
    >>>     bar_patches.sort(key=lambda bar: bar.compute_depth())
    >>>     return bar_patches[-1]
    """
\end{lstlisting}

The example sorts ascending and takes the last element for \emph{furthest away}. The
specification therefore requires that a larger value denote a \emph{greater} distance
from the camera. Section~\ref{sec:sign} shows the shipped implementation violates
this, and that this example is the strongest available evidence for the correction.

\subsection{C2 --- an explicit output contract}

C2 appends three sentences and two examples; nothing is removed.

\begin{lstlisting}[style=plainlist]
This is a VISUAL GROUNDING task. The function must return the single
ImagePatch that the query refers to. It must NOT return a string, a bool,
or a list.
Resolve every spatial relation in the query; do not drop it.
\end{lstlisting}

\begin{lstlisting}
>>> # Query: man on the right
>>> def execute_command(image) -> ImagePatch:
>>>     image_patch = ImagePatch(image)
>>>     man_patches = image_patch.find("man")
>>>     man_patches.sort(key=lambda p: p.horizontal_center)
>>>     return man_patches[-1]

>>> # Query: the cup behind the bottle
>>> def execute_command(image) -> ImagePatch:
>>>     image_patch = ImagePatch(image)
>>>     cup_patches = image_patch.find("cup")
>>>     bottle_patches = image_patch.find("bottle")
>>>     if not cup_patches: return image_patch
>>>     if not bottle_patches: return cup_patches[0]
>>>     bottle_depth = bottle_patches[0].compute_depth()
>>>     behind = [c for c in cup_patches if c.compute_depth() > bottle_depth]
>>>     return behind[0] if behind else cup_patches[0]
\end{lstlisting}

The clause on not dropping spatial relations is included because the static analysis
measured that failure independently of the return type.

\subsection{C3 --- depth primitives}

C3 replaces the first guideline with a distinction between the two kinds of spatial
relation:

\begin{lstlisting}[style=plainlist]
- Use base Python (comparison, sorting) for logical operations and arithmetic.
- Resolve LEFT/RIGHT/ABOVE/BELOW using patch coordinates
  (left, right, horizontal_center, vertical_center).
- Resolve NEAR/FAR/BEHIND/IN FRONT using the depth functions below. These
  relations are about distance from the camera and cannot be recovered from
  image-plane coordinates.
\end{lstlisting}

and documents five new functions:

\begin{lstlisting}[caption={Signatures presented to the generator in C3 and C4.}]
def depth_order(patches: List[ImagePatch],
                reverse: bool = False) -> List[ImagePatch]:
    """Returns the patches sorted from NEAREST to FURTHEST from the camera.
    Pass reverse=True to sort furthest to nearest."""

def is_behind(a: ImagePatch, b: ImagePatch, margin: float = 0.0) -> bool:
    """Returns True if a is further from the camera than b."""

def is_in_front_of(a: ImagePatch, b: ImagePatch, margin: float = 0.0) -> bool:
    """Returns True if a is closer to the camera than b."""

def is_between_3d(a: ImagePatch, b: ImagePatch, c: ImagePatch) -> bool:
    """Returns True if a lies between b and c in depth."""

def distance_3d(a: ImagePatch, b: ImagePatch) -> float:
    """Returns the separation between two patches, combining image-plane
    offset and depth. Relative, not metres."""
\end{lstlisting}

\texttt{depth\_order} is the highest-value addition: a manual audit
(Section~\ref{sec:audit}) finds \emph{closest}, \emph{nearest}, \emph{farthest} and
\emph{closer} account for the majority of depth-word queries, and one ordering
primitive covers all four including the ordinal cases (\emph{second closest}) that a
pairwise predicate cannot express.

\subsection{C4 --- routing}

C4 adds no capability. It states when the capability applies:

\begin{lstlisting}[style=plainlist]
- Use the depth functions ONLY when the query actually expresses a depth
  relation: closest, nearest, farthest, furthest, behind, in front of,
  between (in depth). For every other query they add nothing and cost accuracy.
- If the query is about LEFT/RIGHT/TOP/BOTTOM, sort by patch coordinates and
  do not call any depth function.
- If the query is about a colour, clothing, size or any other attribute, use
  verify_property or best_text_match and do not call any depth function.
\end{lstlisting}

and supplies two counter-examples --- queries that should \emph{not} invoke the depth
API, since the existing examples demonstrate only positive cases:

\begin{lstlisting}
>>> # Query: the man on the left
>>> def execute_command(image) -> ImagePatch:
>>>     image_patch = ImagePatch(image)
>>>     man_patches = image_patch.find("man")
>>>     if not man_patches: return image_patch
>>>     man_patches.sort(key=lambda p: p.horizontal_center)
>>>     return man_patches[0]

>>> # Query: the woman in the red jacket
>>> def execute_command(image) -> ImagePatch:
>>>     image_patch = ImagePatch(image)
>>>     woman_patches = image_patch.find("woman")
>>>     if not woman_patches: return image_patch
>>>     red = [w for w in woman_patches
>>>            if w.verify_property("woman", "red jacket")]
>>>     return red[0] if red else woman_patches[0]
\end{lstlisting}

C4 exists because C3 was measured and found to trade one failure for another; the
sequence is reported in Chapter~\ref{ch:rejected} rather than presented as foresight.

\section{Evaluation protocol}

\subsection{Metric}

The paper's Table~1 is captioned ``we report accuracy on the REC task'' while its
column header reads ``IoU (\%)'', and the released evaluator computes mean IoU. These
are different quantities. This dissertation reports \emph{accuracy at
$\mathrm{IoU}\geq 0.5$} as primary and mean IoU alongside. On the runs reported here
the two differ by roughly one point --- C4 on RefCOCO scores $42.00$ and $41.81$ --- so
no conclusion depends on the choice, but both are given so comparison is possible in
either currency.

\subsection{Handling of invalid predictions}

The released evaluator substitutes the dataset mean box, worth approximately $22.6$
IoU, when a program returns nothing; this implementation scores an invalid return as
zero. The difference is material for C1: under the paper's own evaluator, C1 would
report approximately $22.6$ rather than $0.00$. The programs are equally invalid; what
differs is whether the evaluator conceals the failure. Scoring zero is the more
informative choice for a study whose subject is program validity.

\subsection{Data}

All experiments use testA of RefCOCO and RefCOCO+, with UNC referring expressions and
COCO \texttt{train2014} images. Each condition is evaluated on the first 500
expressions in file order --- deterministic and identical across conditions. This
differs from the published evaluation, which reports on the full split after a seeded
shuffle; the subset is a compute-driven restriction, recorded as such.

Queries are labelled \emph{spatial} or \emph{non-spatial} by a lexicon of relational
terms. The lexicon is a heuristic, not annotation, and Section~\ref{sec:audit}
examines its limitations.

% =====================================================================
\chapter{Implementation}
\label{ch:implementation}

\section{Environment}

Two Conda environments are required and never coexist in one interpreter.

\begin{table}[htbp]\centering
\caption{Execution environments. The GLIP fork requires an older PyTorch than the
generator; the pipeline is split at the \texttt{programs.jsonl} boundary as a result.}
\label{tab:envs}
\begin{tabular}{@{}llll@{}}
\toprule
\textbf{Environment} & \textbf{Stage} & \textbf{PyTorch} & \textbf{Notes} \\
\midrule
\texttt{vipergpt} & generation & 2.13.0 & transformers 5.x \\
\texttt{glip\_env} & execution & 2.1.2+cu121 & transformers 4.36.2, NumPy $<2$ \\
\bottomrule
\end{tabular}
\end{table}

Hardware is a single NVIDIA H100 NVL (95\,GB) allocated through SLURM using fractional
GPU shards. Compute nodes have no outbound network access, so all model weights and
datasets are fetched on the login node and every job asserts offline operation before
loading anything.

\section{Building GLIP for the Hopper architecture}

The published environment pins CUDA 11.6, which has no support for compute capability
9.0. The published setup therefore \emph{cannot run at all} on this hardware. Three
changes were required.

\textbf{Removed headers.} The \texttt{THC/*} headers were deleted from PyTorch after
version 1.11. Seven CUDA source files under \texttt{maskrcnn\_benchmark/csrc/cuda/}
referenced them:

\begin{lstlisting}[style=plainlist]
#include <THC/THCAtomics.cuh>       ->  #include <ATen/cuda/Atomic.cuh>
#include <THC/THCDeviceUtils.cuh>   ->  #include <ATen/cuda/CUDAContext.h>
\end{lstlisting}

\texttt{THCCeilDiv}, \texttt{THCudaMalloc} and \texttt{THCState} were verified absent
from live code, so no further rewriting was needed.

\textbf{Target architectures.} \texttt{setup.py} was extended with explicit
\texttt{sm\_80} and \texttt{sm\_90} targets.

\textbf{NumPy 2.0 aliases.} Five Python files used \texttt{np.float}, removed in NumPy
1.24, causing construction to fail in \texttt{anchor\_generator.py} before any
inference.

The build must be performed \emph{on a GPU node}: \texttt{setup.py} gates
\texttt{CUDAExtension} on \texttt{torch.cuda.is\_available()}, so a login-node build
silently produces a CPU-only extension with no error.

\subsection{Independent verification}

The compiled kernel was verified before the full pipeline was trusted:

\begin{lstlisting}[caption={CUDA NMS verified on a hand-checked case. Boxes 0 and 1
overlap above threshold; box 2 is disjoint.}]
boxes  = torch.tensor([[0.,0.,10.,10.], [1.,1.,11.,11.], [50.,50.,60.,60.]],
                      device='cuda')
scores = torch.tensor([0.9, 0.8, 0.7], device='cuda')
keep   = _C.nms(boxes, scores, 0.5)      # -> [0, 2]
\end{lstlisting}

The detector was then validated against RefCOCO ground truth, reaching a mean IoU of
$0.776$ at threshold $0.2$ and $81.5\%$ top-box accuracy on a held-out sample of 200.

\section{The detector wrapper}

A defect in the first wrapper is worth recording because of how it presented.
\texttt{GLIPDemo} defaults to \texttt{tensor\_inputs=False}, under which it derives the
image size as \texttt{original\_image.shape[:-1]}, assuming an HWC array. Passing a CHW
tensor makes it read $(3, H)$ and rescale every box into a three-pixel-wide image.

\begin{lstlisting}[caption={The one-word correction. Without it, every IoU was
\emph{exactly} zero --- not low, zero --- which is what identified the fault as a
coordinate-space bug rather than a detection-quality problem.}]
self._demo = GLIPDemo(
    glip_cfg,
    min_image_size=800,
    confidence_threshold=self.threshold,
    show_mask_heatmaps=False,
    tensor_inputs=True,      # <- selects shape[-2:] and the tensor transform
)
\end{lstlisting}

Boxes are converted from the detector's top-left origin to the repository's bottom-left
convention on return, and child patches produced by \texttt{find} are offset into their
parent's frame --- an error invisible for root-level programs and wrong for any program
that crops before finding.

\section{The depth backend}
\label{sec:sign}

The shipped wrapper loaded MiDaS through \texttt{torch.hub}, which reaches the network
and therefore hangs on an offline compute node, and returned the raw prediction. MiDaS
predicts \emph{inverse} depth: a larger value means \emph{closer}. Every comparison
built on \texttt{compute\_depth} therefore ran opposite to the semantics the paper's
own example requires.

\begin{lstlisting}[caption={Corrected depth. The inversion lives here and in
\texttt{depth\_of\_region}, and nowhere else; callers are not permitted to read the raw
map, which is what produced three separate sign errors.}]
def compute_depth(self, image) -> float:
    """Distance-like depth of the patch: LARGER = FURTHER from the camera."""
    m = self._inverse_depth_map(image)
    v = self._center_stat(m)
    if v is None:
        return 0.0
    return 1.0 / max(v, _EPS)
\end{lstlisting}

Two further corrections were required. The DPT processor resizes to a square
($384 \times 384$), so the depth map does not share the image's aspect ratio;
proportional region reads therefore sampled horizontally squashed regions. The map is
now resampled to image size. And the depth statistic is taken over the central 50\% of
a box, because border pixels of a detection are usually background and skew the median.

\section{The depth primitives}

\begin{lstlisting}[caption={The contributed primitives, from
\texttt{pipeline/image\_patch.py}. All follow the convention LARGER = FURTHER.}]
def depth_order(patches: list, reverse: bool = False) -> list:
    """Patches sorted NEAR to FAR. reverse=True gives far to near.

    Backs "closest", "nearest", "farthest", "second closest" and
    "third from the front" with a single primitive. Patches whose depth
    cannot be computed are dropped rather than silently placed at one end.
    """
    if not patches:
        return []
    scored = []
    for p in patches:
        try:
            d = p.compute_depth()
        except Exception:          # a bad crop must not kill the program
            continue
        if d is not None:
            scored.append((p, d))
    scored.sort(key=lambda t: t[1], reverse=reverse)
    return [p for p, _ in scored]


def is_behind(a, b, margin: float = 0.0) -> bool:
    """True if a is further from the camera than b by at least margin."""
    return a.compute_depth() > b.compute_depth() + margin


def is_between_3d(a, b, c) -> bool:
    """True if a lies between b and c in depth."""
    lo, hi = sorted([b.compute_depth(), c.compute_depth()])
    return lo < a.compute_depth() < hi
\end{lstlisting}

The primitives are exposed to generated programs by the executor's namespace builder,
alongside the original API.

\subsection{Unit tests}

Three sign errors occurred during development, none of which was caught by a test.
The ordering semantics are now pinned by eighteen assertions over stub patches with
known depths, requiring no model and no GPU:

\begin{lstlisting}[caption={Representative assertions from
\texttt{test\_depth\_api.py}. A failure here means the sign or the ordering has
flipped, and is visible immediately rather than after hours of debugging.}]
near, mid, far = StubPatch("near", 1.0), StubPatch("mid", 5.0), StubPatch("far", 9.0)

assert depth_order([far, near, mid])            == [near, mid, far]
assert depth_order([near, far, mid], True)      == [far, mid, near]
assert depth_order([far, mid, near])[0]         is near   # "closest"
assert depth_order([near, mid, far])[-1]        is far    # "farthest"
assert depth_order([far, StubPatch("bad", None), near]) == [near, far]
assert is_behind(far, near) and not is_behind(near, far)
assert is_between_3d(mid, near, far)
assert round(distance_3d(a, b), 6) == 5.0                 # 3-4-5 triangle
\end{lstlisting}

\section{Execution and scoring}

The execution engine runs each program in a restricted namespace with a wall-clock
timeout, and returns errors as data rather than raising, so that one bad program does
not terminate a 500-sample run.

\begin{lstlisting}[caption={Scoring, from \texttt{eval/execute\_refcoco.py}. A non-patch
return scores zero rather than being excluded --- returning a sentence where a box is
required is a task failure, and dropping such cases would flatter the baseline.}]
res  = execute_program(p["program"], chw, bus, timeout_s=args.timeout)
kind = "error" if res.error else classify(res.result)

iou, pred = 0.0, None
if kind == "patch":
    r = res.result
    # ImagePatch is bottom-left origin; undo the flip for COCO comparison
    pred = (r.left, H - r.upper, r.right, H - r.lower)
    iou  = iou_xyxy(pred, g["gt_xyxy"])
ious.append(iou)          # non-patch returns score 0 -- that is the point
\end{lstlisting}

Because the programs are generated in one environment and executed in another, the two
stages are joined by index. The runner asserts that the query string of every program
matches the query of the corresponding ground-truth record, and refuses to score if any
disagree.

% =====================================================================
\chapter{Experimental Setup}
\label{ch:setup}

\begin{table}[htbp]\centering
\caption{Experiments. Generation and execution are separate jobs; identifiers are
scheduler references retained for provenance.}
\label{tab:experiments}
\begin{tabular}{@{}llll@{}}
\toprule
\textbf{Experiment} & \textbf{Varies} & \textbf{Holds fixed} & \textbf{Jobs} \\
\midrule
Static analysis & prompt & model, data & 29187, 29193 \\
End-to-end accuracy & prompt & model, modules, data & 29290, 29292 \\
Generator scale & model size & prompt, modules, data & 29321, 29352, 29404 \\
Seed variance & decoding seed & everything else & 29301, 29306, 29526 \\
Depth conditions & prompt (C3, C4) & model, modules, data & 29459, 29523--29525 \\
Threshold sweep & detector threshold & held-out sample & 29318 \\
Candidate cap & candidate-set size & programs & 29563 \\
Cap control & cap $\times$ condition & programs & 29565 \\
\bottomrule
\end{tabular}
\end{table}

Decoding is greedy for all headline figures, matching the published temperature of
zero. Because greedy decoding is deterministic, variance is estimated separately by
sampling at $T = 0.7$ with three seeds; those figures are reported as mean and standard
deviation and are never mixed with the greedy values.

\section{Detector threshold}

The paper specifies a detection threshold of $0.5$. That value was found to lose eight
of twenty detections outright, so a sweep was performed --- on RefCOCO/testA samples
200--399, \emph{disjoint} from every sample on which results are reported, so that the
threshold is not tuned on the data it is evaluated on.

\begin{table}[htbp]\centering
\caption{Detector threshold sweep, held-out sample of 200. The selected value of $0.2$
lies within one point of the optimum for detection.}
\label{tab:threshold}
\begin{tabular}{@{}rrrrr@{}}
\toprule
\textbf{Threshold} & \textbf{Acc.\ IoU$\geq$0.5} & \textbf{Mean IoU} &
\textbf{No detection} & \textbf{Candidates/image} \\
\midrule
0.10 & 81.50 & 0.7372 & 21 & 25.5 \\
0.15 & 81.50 & 0.7372 & 21 & 25.4 \\
\textbf{0.20} & \textbf{80.50} & 0.7336 & 21 & \textbf{19.7} \\
0.25 & 70.00 & 0.6658 & 21 & 9.3 \\
0.30 & 59.50 & 0.5784 & 22 & 5.4 \\
0.40 & 40.00 & 0.4350 & 32 & 2.4 \\
0.50 & 31.00 & 0.3360 & 56 & 1.3 \\
\bottomrule
\end{tabular}
\end{table}

\begin{figure}[htbp]\centering
\includegraphics[width=0.66\textwidth]{results/figures/fig8_threshold_tradeoff.pdf}
\caption[Detector threshold trade-off]{One threshold, two incompatible jobs. The value
that maximises top-box detection accuracy also returns roughly twenty candidates per
image; a program that must \emph{select} among candidates wants far fewer. The same
knob cannot serve both, a tension that recurs in Chapter~\ref{ch:rejected}.}
\label{fig:threshold}
\end{figure}

Twenty-one of two hundred samples yield no detection at \emph{every} threshold up to
$0.4$. That floor is not threshold-related: it arises from the head-noun heuristic used
to prompt the detector, under which \emph{``player number 8''} reduces to \emph{``8''}
and is discarded.

% =====================================================================
\chapter{Results}
\label{ch:results}

\section{Main result}

\begin{table}[htbp]\centering
\caption{Accuracy at IoU $\geq 0.5$. Qwen2.5-Coder-7B, greedy decoding, $n=500$ per
cell, testA split. C1 is the released prompt; C4 adds both the depth primitives and
the rule stating when they apply.}
\label{tab:main}
\begin{tabular}{@{}llrrrr@{}}
\toprule
\textbf{Dataset} & \textbf{Subset} & \textbf{C1} & \textbf{C2} & \textbf{C3} & \textbf{C4} \\
\midrule
RefCOCO+ & overall & 0.00 & 30.80 & 28.00 & \up{37.60} \\
RefCOCO+ & spatial ($n{=}47$) & 0.00 & 4.26 & 19.15 & \up{21.28} \\
RefCOCO+ & non-spatial ($n{=}453$) & 0.00 & 33.55 & \down{28.92} & \up{39.29} \\
\midrule
RefCOCO & overall & 0.00 & 39.20 & 36.40 & \up{42.00} \\
RefCOCO & spatial ($n{=}289$) & 0.00 & 34.26 & 36.68 & \up{37.02} \\
RefCOCO & non-spatial ($n{=}211$) & 0.00 & 45.97 & \down{36.02} & \up{48.82} \\
\bottomrule
\end{tabular}
\end{table}

\begin{figure}[htbp]\centering
\includegraphics[width=\textwidth]{results/figures/fig1_main_result.pdf}
\caption[Main result]{All four conditions across both datasets and both subsets. The
three rightmost groups act as controls: an intervention that improved only its target
subset would be a trade rather than a gain.}
\label{fig:main}
\end{figure}

Mean IoU for the same RefCOCO runs is $0.0000$, $0.3875$, $0.3655$ and $0.4181$ for
C1--C4, tracking accuracy to within about a point throughout.

\section{Comparison with the published figures}

\begin{figure}[htbp]\centering
\includegraphics[width=\textwidth]{results/figures/fig16_paper_vs_result.pdf}
\caption[Reproduction against published figures]{Reproduced accuracy against the
paper's Table~1, with the zero-shot baselines it reports alongside. C4 recovers 58\% of
the RefCOCO figure and 56\% of RefCOCO+.}
\label{fig:vspaper16}
\end{figure}

\begin{figure}[htbp]\centering
\includegraphics[width=\textwidth]{results/figures/fig10_vs_paper.pdf}
\caption[Reproduction with the perception ceiling]{The dotted line is GLIP alone taking
its single best box: $81.5\%$ on a held-out sample of 200 --- \emph{above} the paper's
$72.0$. Perception is therefore not what separates this reproduction from the original.
No decomposition of the gap into individual deviations is shown, because those
ablations were not run.}
\label{fig:vspaper}
\end{figure}

That the perception ceiling exceeds the published figure is informative. The detector
available here is good enough to surpass $72.0$ if a program selected from it perfectly.
The shortfall is therefore in program synthesis, which is measured support for treating
the generator substitution (D1) as the dominant term rather than assuming it.

\section{Why the released prompt scores zero}
\label{sec:static}

The result is not a crash. Of 500 programs generated under C1 on RefCOCO, 332 executed
to completion and returned a \emph{string}; none returned an \texttt{ImagePatch}.

\begin{table}[htbp]\centering
\caption{What each of the 500 samples returned, RefCOCO/testA.}
\label{tab:returnkinds}
\begin{tabular}{@{}lrrr@{}}
\toprule
\textbf{Return type} & \textbf{C1} & \textbf{C2} & \textbf{C4} \\
\midrule
\texttt{ImagePatch} & 0 & 444 & 488 \\
\texttt{str} & 332 & 1 & 0 \\
error & 168 & 55 & 11 \\
\bottomrule
\end{tabular}
\end{table}

\begin{figure}[htbp]\centering
\includegraphics[width=\textwidth]{results/figures/fig7_failure_modes.pdf}
\caption[Failure decomposition]{Every sample classified by what happened to it. C1
fails on \emph{return type} and never reaches localisation; by C4 those categories are
gone and the residual failure is a localisation miss, which is a perception limit
rather than a specification bug. An accuracy figure alone conflates the two.}
\label{fig:failure}
\end{figure}

Missing modules are ruled out as a cause. \texttt{simple\_query} and
\texttt{llm\_query} are given neutral empty-string fallbacks (D10), so no program can
fail for lack of a module that was deliberately not loaded. That concession strictly
favours C1 and does not change its score.

\subsection{A second effect: output format}

Under C1 a large fraction of generations never defined \texttt{execute\_command} at
all --- the model emitted prose or malformed markdown: $140/500$ on RefCOCO and
$212/500$ on RefCOCO+, against zero under every later condition. The contract
disciplines output \emph{format} as well as return type; these are separable failures
and are reported separately.

\subsection{Static analysis}

Before execution, the programs were analysed by parsing.

\begin{table}[htbp]\centering
\caption{Static analysis of generated programs, $n=500$ per cell. ``Task-invalid''
denotes a program whose \texttt{return} yields a string rather than a patch;
``spatial dropped'' denotes a spatial query whose program contains no positional
operation at all.}
\label{tab:static}
\begin{tabular}{@{}lrrrr@{}}
\toprule
\textbf{Metric} & \textbf{C1 RefCOCO} & \textbf{C2 RefCOCO} & \textbf{C1 RefCOCO+} & \textbf{C2 RefCOCO+} \\
\midrule
Task-invalid return & 98.8\% & \up{0.2\%} & 98.2\% & \up{0.0\%} \\
Spatial constraint dropped & 17.3\% & \up{1.7\%} & 31.9\% & \up{2.1\%} \\
Parse rate & 99.8\% & \down{96.6\%} & 99.6\% & \down{92.2\%} \\
Mean API calls per program & 2.45 & 1.66 & 2.89 & 2.25 \\
\bottomrule
\end{tabular}
\end{table}

The contract costs parse rate --- $3.2$ points on RefCOCO and $7.4$ on RefCOCO+ ---
because a stricter requirement induces a small rate of syntactically invalid
generations. Net structural validity nonetheless rises from roughly $1\%$ to $96\%$.

\section{Generator scale}

\begin{table}[htbp]\centering
\caption{Generator scale, condition C2, greedy, RefCOCO/testA, $n=500$.}
\label{tab:scale}
\begin{tabular}{@{}lrrrrr@{}}
\toprule
\textbf{Generator} & \textbf{C1 acc.} & \textbf{C2 acc.} & \textbf{spatial} &
\textbf{non-spatial} & \textbf{C1 task-invalid} \\
\midrule
Qwen2.5-Coder-1.5B & 0.80 & 34.40 & 34.95 & 33.65 & 95.8\% \\
Qwen2.5-Coder-7B & 0.00 & 39.20 & 34.26 & 45.97 & 98.8\% \\
Qwen2.5-Coder-32B & 0.00 & 40.80 & 34.95 & 48.82 & 87.6\% \\
\bottomrule
\end{tabular}
\end{table}

\begin{figure}[htbp]\centering
\includegraphics[width=0.62\textwidth]{results/figures/fig2_scale.pdf}
\caption[Accuracy against generator scale]{Non-spatial accuracy climbs with capacity;
spatial accuracy does not move. Same images, same detector, same prompt --- only the
generator differs. This is the evidence that the spatial deficit is an interface
problem rather than a capacity one.}
\label{fig:scale}
\end{figure}

Two conclusions follow. The specification gap is scale-independent: a 32B generator
still emits $87.6\%$ task-invalid programs under C1 and still scores $0.00$, and the
same three sentences repair it at every size. Scaling $1.5\mathrm{B} \to 32\mathrm{B}$
buys $6.4$ accuracy points; the contract buys $33$ to $41$.

More consequentially, \textbf{spatial accuracy does not respond to scale at all} ---
$34.95 / 34.26 / 34.95$, with the smallest and largest models identical to two decimal
places --- while non-spatial accuracy rises by $15.2$ points. A larger generator writes
better programs for everything except spatial relations, where it writes programs that
are exactly as wrong. If the deficit were a program-synthesis problem, capacity would
help; it demonstrably does elsewhere.

An incidental observation: under C1, the 32B model failed to define
\texttt{execute\_command} in $304/500$ generations against $16/500$ for 1.5B. The
larger instruction-tuned model is more prone to answering in prose when the prompt does
not demand a function, which is why 32B scores $0.00$ where 1.5B manages $0.80$.

\section{Spatial breakdown}
\label{sec:audit}

We expected the spatial gap to \emph{shrink} on RefCOCO+, which forbids location words.
It widened threefold: $9.3$ points on RefCOCO against $28.1$ on RefCOCO+ under C2.

A manual audit of all 47 RefCOCO+ spatial queries explains why. The two subsets contain
different kinds of spatial language.

\begin{table}[htbp]\centering
\caption{The two spatial subsets are not the same kind of query.}
\label{tab:audit}
\begin{tabular}{@{}p{0.16\textwidth}p{0.34\textwidth}p{0.38\textwidth}@{}}
\toprule
& \textbf{RefCOCO} & \textbf{RefCOCO+} \\
\midrule
Type & 2D image-plane & depth and inter-object \\
Examples & \emph{man on right}, \emph{person bottom left} &
\emph{man closest to us}, \emph{farthest man}, \emph{umpire behind catcher} \\
Resolvable by & sorting patch coordinates in base Python &
requires \texttt{compute\_depth} or 3D geometry \\
Accuracy (C2) & 34.26 & 4.26 \\
\bottomrule
\end{tabular}
\end{table}

Forbidding location words removed viewer-frame \emph{directional} expressions and left
depth-ordering and relational ones. The released prompt's instruction to use base Python
for ``left/right/up/down'' works for the first and cannot work for the second.
Consistently, the call-frequency analysis finds \texttt{compute\_depth} invoked 29 times
against 547 for \texttt{find}: the geometric API is effectively unused, by design of the
prompt.

Two further observations support this reading. All three successes in the RefCOCO+
spatial subset under C2 are depth queries with very high IoU ($0.97$, $0.96$, $0.92$):
when depth ordering is attempted correctly the localisation is excellent, but it is
usually not attempted. And across seeds, spatial variance is roughly four times
non-spatial variance, reported in Section~\ref{sec:seeds}.

\textbf{Caveat.} $n = 47$, so $4.26\%$ is 2 of 47 and $21.28\%$ is 10 of 47. The lexicon
also produces roughly five false positives in that subset --- \emph{``writing on back''},
\emph{``long sleeved''} matching \emph{side} --- which would push true depth-query
accuracy lower still. The direction is established; the value is not.

\section{API adoption}

The principal risk in extending an API is that the generator ignores the extension, as
it largely ignores \texttt{compute\_depth}. It did not.

\begin{table}[htbp]\centering
\caption{API calls across 500 generated programs, RefCOCO/testA.}
\label{tab:adoption}
\begin{tabular}{@{}lrrrr@{}}
\toprule
\textbf{Condition} & \texttt{find} & \texttt{compute\_depth} & \texttt{depth\_order} &
\textbf{other depth primitives} \\
\midrule
C1 & 548 & 29 & 0 & 0 \\
C2 & 639 & 23 & 0 & 0 \\
C3 & 598 & 19 & 119 & 63 \\
C4 & 559 & --- & 37 & 7 \\
\bottomrule
\end{tabular}
\end{table}

\begin{figure}[htbp]\centering
\includegraphics[width=0.68\textwidth]{results/figures/fig3_api_adoption.pdf}
\caption[API adoption across conditions]{Zero adoption under the older prompts, strong
adoption once the primitives are offered. Under C4 the count falls to 37 --- the routing
rule working, not the primitive being abandoned, since RefCOCO contains few depth
queries. On RefCOCO+, which has 47 depth-word queries, C4 makes 33
\texttt{depth\_order} calls against 167 under C3.}
\label{fig:adoption}
\end{figure}

\section{Providing a capability is not sufficient}

C3 offers the primitives; C4 adds three sentences saying when they apply. The difference
between them is the sharper finding of this dissertation.

C3 wins its target subset --- RefCOCO+ depth queries move $4.26 \to 19.15$ --- and loses
elsewhere: non-spatial accuracy falls $45.97 \to 36.02$ on RefCOCO and $33.55 \to 28.92$
on RefCOCO+, so overall accuracy drops. \texttt{depth\_order} is invoked 119 times on
RefCOCO, where most queries need no depth at all. \emph{Given a new primitive, the
generator reaches for it indiscriminately.}

C4 improves every cell of both datasets. The depth gain survives ($21.28$) while
non-spatial rises above C2 as well ($48.82$ against $45.97$), presumably because the
counter-examples also clarify the 2D and attribute paths. Execution robustness improves
monotonically: RefCOCO errors $55 \to 24 \to 11$ and patch returns $444 \to 475 \to 488$
for C2 $\to$ C3 $\to$ C4.

This is the same class of omission as the first result. In C1 the specification failed to
state \emph{what} to return; in C3 it failed to state \emph{when} a capability applies.
Both are failures of the contract rather than of the model.

\begin{figure}[htbp]\centering
\includegraphics[width=\textwidth]{results/figures/fig4_qualitative.pdf}
\caption[Qualitative comparison]{Depth queries where C4 succeeds and C2 fails. Red: C2's
prediction. Teal: C4's. Dashed: ground truth. The failures are not near-misses --- without
a primitive for \emph{closest}, the program selects a different object entirely.}
\label{fig:qual}
\end{figure}

\section{Seed variance}
\label{sec:seeds}

Greedy decoding is deterministic, so variance was estimated by sampling at $T = 0.7$ with
three seeds.

\begin{table}[htbp]\centering
\caption{Seed variance, RefCOCO/testA, $T=0.7$, seeds 1--3, $n=500$.}
\label{tab:seeds}
\begin{tabular}{@{}lccc@{}}
\toprule
\textbf{Condition} & \textbf{overall} & \textbf{spatial} & \textbf{non-spatial} \\
\midrule
C2 contract & $36.07 \pm 1.75$ & $32.87 \pm 3.35$ & $40.44 \pm 0.72$ \\
C3 + depth & $34.40 \pm 0.40$ & $34.02 \pm 0.80$ & $34.92 \pm 0.55$ \\
C4 + routing & $\mathbf{42.07 \pm 0.61}$ & $38.29 \pm 0.72$ & $47.23 \pm 1.19$ \\
\bottomrule
\end{tabular}
\end{table}

\begin{figure}[htbp]\centering
\includegraphics[width=\textwidth]{results/figures/fig6_seed_variance.pdf}
\caption[Seed variance and stability]{Left: C4 separates from C2 by $6.0$ points against
standard deviations under two. Right: standard deviation on spatial queries alone. The
explicit primitive removes the improvisation the generator otherwise performs afresh on
every sample.}
\label{fig:seeds}
\end{figure}

Two observations. First, \textbf{greedy decoding understates C4}. C2's greedy result
($39.20$) sits $1.8$ standard deviations \emph{above} its own sampled mean ($36.07$) ---
greedy was fortunate for C2 --- while C4's greedy ($42.00$) and sampled ($42.07$) results
coincide. The honest comparison is $6.0$ points, not the $2.8$ the greedy table suggests.

Second, and unanticipated: \textbf{the primitives buy stability as well as accuracy}.
Spatial-query standard deviation falls from $\pm 3.35$ under C2 to $\pm 0.72$ under C4, a
factor of $4.6$. Under C2 the generator improvises spatial logic on every sample ---
sometimes sorting by coordinate, sometimes dropping the relation entirely --- and which it
does depends on the seed. Once \texttt{depth\_order} exists there is one way to express
the relation and the variance collapses. An explicit primitive makes the system's
behaviour reproducible, not merely more accurate.

\section{Robustness of the metric}

\begin{figure}[htbp]\centering
\includegraphics[width=\textwidth]{results/figures/fig11_iou_sweep.pdf}
\caption[IoU threshold sensitivity]{Accuracy recomputed at every IoU threshold from $0.3$
to $0.9$, and the C4$-$C2 difference. The threshold of $0.5$ is used because the paper
uses it, but the C4 advantage stays positive across the whole range.}
\label{fig:iousweep}
\end{figure}

Were the advantage confined to the neighbourhood of $0.5$, it would suggest the conditions
differ in box tightness rather than in object selection. They do not.

\begin{figure}[htbp]\centering
\includegraphics[width=\textwidth]{results/figures/fig5_iou_distribution.pdf}
\caption[IoU distribution]{Left: distribution over samples scoring anything at all.
Right: share of samples scoring exactly zero. The distributions are strongly bimodal ---
programs either localise well or miss entirely --- which is what one expects when the
failure mode is selecting the wrong object rather than drawing a loose box.}
\label{fig:distribution}
\end{figure}

% =====================================================================
\chapter{Rejected Hypotheses and Controls}
\label{ch:rejected}

Three hypotheses were tested and rejected, and one intervention was disqualified by a
control. They are reported because a reproduction study in which nothing failed would be
less credible, and because in two cases the rejected result was larger than the accepted
one.

\section{Rejection 1: the collapse hypothesis}

Milestone~1 was designed to test a specific claim: that spatial-relational queries would
cause the generator to collapse the whole query into a single opaque
\texttt{simple\_query} call. The paper's GQA prompt instructs exactly this, listing
relational terms and directing the model to ask directly rather than decompose.

The measured collapse rate was approximately $0\%$ --- one program in 289 spatial queries.

The hypothesis was well-posed and aimed at the wrong task. That instruction applies to
GQA, not to grounding, and grounding cannot use \texttt{simple\_query} at all because it
must return a patch. The cost of establishing this was a single 63-minute job, incurred
before any perception model had been built --- which is why the generation stage was
studied first.

The same run produced the return-type finding that became the first result of this
dissertation.

\section{Rejection 2: RefCOCO+ as a clean control}

RefCOCO+ forbids location words, so it was expected to serve as a control in which the
spatial/non-spatial gap would narrow or vanish. It widened threefold
(Section~\ref{sec:audit}).

The audit that followed produced the distinction between image-plane and depth-based
spatial language, which is the entire justification for Milestone~3. A rejected
hypothesis produced the contribution.

\section{Rejection 3: the primitives alone}

C3 was expected to suffice. It improved its target subset and degraded the others,
failing criterion C-2, and was superseded by C4 (Section~\ref{sec:prompts}).

\section{A fifteen-point result, disqualified}

The detector returns roughly 20 candidates per image at the selected threshold. Capping
\texttt{find} to $k$ of them, on the \emph{same} C4 programs:

\begin{table}[htbp]\centering
\caption{Candidate cap applied to C4 programs, RefCOCO/testA.}
\label{tab:cap}
\begin{tabular}{@{}lrrr@{}}
\toprule
\textbf{Cap} & \textbf{overall} & \textbf{spatial} & \textbf{non-spatial} \\
\midrule
off & 42.80 & 37.72 & 49.76 \\
$k = 12$ & 45.20 & 41.52 & 50.24 \\
$k = 6$ & 49.60 & 47.75 & 52.13 \\
$k = 3$ & \textbf{57.60} & \textbf{61.59} & 52.13 \\
\bottomrule
\end{tabular}
\end{table}

$57.60$ would be the best figure in this project by fifteen points, and $61.59$ on spatial
queries would exceed every non-spatial result obtained. It is not reportable, and the
control is why.

Applying the identical cap to \textbf{C2}, which contains no depth primitives at all:

\begin{table}[htbp]\centering
\caption{The control required by criterion C-1. C2 has no depth reasoning anywhere in the
pipeline.}
\label{tab:capcontrol}
\begin{tabular}{@{}lrrr@{}}
\toprule
& \textbf{no cap} & \textbf{$k = 3$} & \textbf{gain} \\
\midrule
C2 spatial & 34.60 & 57.09 & \down{$+22.5$} \\
C4 spatial & 37.72 & 61.59 & $+23.9$ \\
C2 non-spatial & 45.02 & 45.50 & $+0.5$ \\
C4 non-spatial & 49.76 & 52.13 & $+2.4$ \\
\bottomrule
\end{tabular}
\end{table}

\begin{figure}[htbp]\centering
\includegraphics[width=\textwidth]{results/figures/fig13_cap_control.pdf}
\caption[The control that disqualified a fifteen-point result]{Left: both conditions with
and without a three-candidate cap. Right: the gain the cap alone provides, by subset. C2
gains almost exactly what C4 gains, so the improvement belongs to the candidate set rather
than to depth reasoning. Non-spatial queries barely move, which identifies the mechanism.}
\label{fig:capcontrol}
\end{figure}

C2's spatial accuracy rises $22.5$ points from the cap alone. Of C4's $23.9$, roughly
$1.4$ points are attributable to depth and the rest to the candidate set.

\textbf{The mechanism is arithmetic.} Choosing among three candidates rather than twenty
raises the hit rate by construction --- random selection alone scores $33\%$ at $k=3$
against $5\%$ at $k=20$ --- and in RefCOCO the referred object is almost always among the
most confident detections. That non-spatial queries gain almost nothing ($+0.5$ for C2)
confirms it: those queries verify attributes rather than selecting among same-class
candidates, so a smaller set does not help them. Only \emph{pick one of $N$ people}
benefits, and it benefits however the pick is made.

\textbf{Decision.} \texttt{max\_detections} remains disabled for every reported number.
Non-maximum suppression is likewise left off: it helps RefCOCO marginally
($42.00 \to 42.80$) but harms RefCOCO+ depth queries badly ($21.28 \to 12.77$), because
suppressing overlapping boxes removes exactly the depth-distinct instances a depth query
must tell apart.

\textbf{What survives.} C4's advantage over C2 holds under both settings --- $3.8$ points
uncapped, $5.4$ capped --- so no claim in this work rests on the cap.

\begin{figure}[htbp]\centering
\includegraphics[width=0.68\textwidth]{results/figures/fig12_headroom.pdf}
\caption[Achieved against achievable]{Depth-ordering selection study. Grey: the correct box
was present among the candidates. Teal: it was selected. The gap between bars is an
ordering failure; the gap from grey to $100\%$ is a detection failure. The binding
constraint has moved from ordering to detection.}
\label{fig:headroom}
\end{figure}

\section{A sign error that a synthetic test endorsed}

The depth sign was established, reversed, and re-established during development. The
reason is worth recording.

The initial implementation returned the raw MiDaS prediction. Believing this to be inverse
depth, an inversion was applied. A synthetic test scene --- a brightness gradient with a
large dark object low in frame --- then indicated the inversion was wrong, and it was
removed. That scene was invalid: the model read the brightness gradient as depth.

The decisive experiment removes both the detector and the synthetic data. Take RefCOCO+
depth queries, use \textbf{ground-truth boxes only}, and ask whether the annotated target
sits at the correct end of the depth ordering.

\begin{table}[htbp]\centering
\caption{Depth sign validated against ground truth, no detector involved, $n=25$.}
\label{tab:sign}
\begin{tabular}{@{}lr@{}}
\toprule
& \textbf{Target at correct depth extreme} \\
\midrule
without inversion & \down{$2/25 = 8\%$} \\
with inversion & \up{$20/25 = 80\%$} \\
chance ($3.8$ objects per image) & $26\%$ \\
\bottomrule
\end{tabular}
\end{table}

\begin{figure}[htbp]\centering
\includegraphics[width=0.56\textwidth]{results/figures/fig9_depth_sign.pdf}
\caption[Depth sign validation]{$8\%$ is far \emph{below} the chance level of $26\%$. A
weak signal scores near chance; an inverted one scores below it. That distinction is what
identified the fault.}
\label{fig:sign}
\end{figure}

Two lessons are recorded. \textbf{Synthetic scenes are unreliable for validating
perceptual signals}, because the model latches onto whatever cue the synthetic image
actually contains. And the paper's own \texttt{compute\_depth} example ---
\texttt{sort(\ldots); return bar\_patches[-1]} for \emph{furthest away} --- independently
requires the corrected sign, so the correction is not merely empirically supported but
required by the specification being reproduced.

% =====================================================================
\chapter{Deviations from the Published Configuration}
\label{ch:deviations}

Every departure is recorded with its reason and expected effect, per criterion C-5.

\begin{longtable}{@{}p{0.05\textwidth}p{0.26\textwidth}p{0.60\textwidth}@{}}
\caption{Deviation register.}\label{tab:deviations}\\
\toprule
\textbf{ID} & \textbf{Deviation} & \textbf{Reason and expected effect} \\
\midrule
\endfirsthead
\toprule
\textbf{ID} & \textbf{Deviation} & \textbf{Reason and expected effect} \\
\midrule
\endhead
\bottomrule
\endfoot

D1 & Codex $\to$ Qwen2.5-Coder & Codex is retired; no faithful option exists. This
substitution is the subject of the study and the dominant term in the gap to $72.0$. \\

D2 & Scope restricted to grounding & The four-table configuration requires roughly
250\,GB against 231\,GB of quota; NExT-QA alone needs over 100\,GB of video, and OK-VQA
depends on GPT-3, also retired. No cross-task claim is made. \\

D3 & RefCOCO annotations from the Internet Archive & The canonical host no longer
resolves. The archived files are byte-identical. No effect. \\

D4 & Dependency versions & The environment pre-existed, so the pinned versions were not
applied. Frozen to a lockfile. No effect on the text-only path. \\

D5 & GLIP built for \texttt{sm\_90} & CUDA 11.6 has no support for compute capability
9.0, so the published environment cannot run on this hardware. Headers updated, target
architectures added; kernels unchanged. Verified numerically. \\

D6 & NumPy 2.0 aliases in GLIP & Five files used \texttt{np.float}, removed in NumPy
1.24. Semantics unchanged. \\

D7 & Offline operation & GLIP fetches a tokenizer at load time; compute nodes have no
network. Cached in advance and asserted before every job. \\

D8 & Detector threshold $0.5 \to 0.2$ & The paper's value loses 8 of 20 detections.
Re-swept on a held-out slice ($n{=}200$, disjoint from all reported cells): $0.2$ is
within $1.0$ point of optimal for detection. Material and now validated. \\

D9 & X-VLM $\to$ CLIP ViT-L & X-VLM weights are distributed via a file-sharing link and
cannot be fetched by script. CLIP scores globally rather than at region level, so
attribute verification should be weaker; costliest on RefCOCO+, whose expressions are
attribute-based. \\

D10 & \texttt{simple\_query}, \texttt{llm\_query} return \texttt{""} & Loading BLIP-2-XXL
and an 8B text model is out of scope. Without a fallback, a program calling them raises,
making its failure attributable to our configuration. Strictly favours C1; no effect on
later conditions. \\

D11 & Depth backend corrected & The shipped wrapper loaded MiDaS over the network and
returned inverse depth unmodified, so comparisons ran opposite to the documented
semantics. Now loaded offline, resampled to image size, and inverted. Every depth number
depends on this; pinned by 18 unit assertions. \\

D12 & Detection scores exposed; candidate controls added & \texttt{find} now also returns
per-box confidence, with optional \texttt{max\_detections} and NMS. \textbf{Both disabled
by default}: enabling them raises accuracy by 15 points, but a control shows the gain
belongs to the candidate set rather than to the method
(Chapter~\ref{ch:rejected}). Recorded as a negative result. \\
\end{longtable}

\section{Undocumented at the time of writing}

Two further differences were identified during the audit of this dissertation and are
recorded here for completeness.

The \texttt{crop\_larger\_margin} option, which expands detections by 10\%, is enabled in
the published configuration and disabled in the evaluation configurations used here. Its
effect has not been measured.

The published evaluator shuffles the split under a fixed seed and reports on the full
testA set; this work takes the first 500 expressions in file order. Both are
deterministic, but they are not the same 500 samples.

% =====================================================================
\chapter{Limitations}
\label{ch:limitations}

\begin{enumerate}[leftmargin=1.6em, itemsep=4pt]
\item \textbf{The reproduction is not faithful.} $42.00$ against a reported $72.0$. The
contributors are enumerated in Chapter~\ref{ch:deviations} but not decomposed
quantitatively, because the per-deviation ablations were not run.

\item \textbf{Metric and fallback differ from the published evaluator.} Accuracy at
IoU $\geq 0.5$ rather than mean IoU, and zero rather than a mean-box substitute for an
invalid prediction. Both figures are reported wherever a comparison is drawn.

\item \textbf{The depth subset is small.} $n = 47$; $21.28\%$ is 10 of 47. The direction
is established, the value is not. The lexicon used to identify spatial queries is a
heuristic with an estimated five false positives in that subset.

\item \textbf{Seed coverage is asymmetric.} Seeds were run for RefCOCO only; every
RefCOCO+ figure, including the headline $21.28$, is single-run.

\item \textbf{X-VLM was never built.} CLIP is a substitute whose weakness is expected to
fall most heavily on RefCOCO+.

\item \textbf{Evaluation is on a 500-sample subset} of testA rather than the full split.

\item \textbf{The candidate-set tension is unresolved.} The detector threshold is correct
for detection and wrong for selection; \texttt{max\_detections} addresses this but was
disqualified as a contribution and disabled.

\item \textbf{Provenance is by job identifier rather than commit hash} for the execution
runs, because \texttt{git} is not installed on the compute nodes. The job scripts are
committed, so the configuration is recoverable, but the automatic stamping intended by
the design is absent.
\end{enumerate}

% =====================================================================
\chapter{Conclusion and Future Work}
\label{ch:conclusion}

\section{Conclusion}

This dissertation set out to reproduce ViperGPT after the withdrawal of its program
generator, and to determine which of its reported behaviour belonged to the method rather
than to that generator.

The answer to RQ1 is that the released specification carries an unstated output contract.
With an open-weights generator, $98.8\%$ of programs return a natural-language string
where a bounding box is required, and executed accuracy is $0.00$ --- not because the
system crashes, but because $332$ of $500$ programs run to completion and answer the wrong
kind of question. Three sentences stating the contract raise accuracy to $39.20$.

The answer to RQ2 is that this is a property of the specification, not of capacity. The
failure occurs at 1.5B, 7B and 32B parameters and is repaired at each. Scaling the
generator twentyfold buys approximately six accuracy points where stating the contract
buys thirty-three to forty-one.

The answer to RQ3 is that accuracy differs sharply by the kind of spatial language. Across
the same three model sizes, non-spatial accuracy rises $33.65 \to 48.82$ while spatial
accuracy remains flat at $34.95 / 34.26 / 34.95$. A manual audit shows that RefCOCO's
spatial queries are image-plane and RefCOCO+'s are depth-based, and that the latter score
$4.26$ against the former's $34.26$. The API exposes one scalar for depth and the prompt
directs the generator away from geometry; the deficit is in the interface.

The answer to RQ4 is qualified and is the more interesting result. Extending the API with
depth-ordering primitives improves depth queries from $4.26$ to $19.15$ but costs ten
points of non-spatial accuracy: the generator applies a new primitive indiscriminately.
Adding three sentences that delimit when the primitives apply improves \emph{every} subset
of both datasets, reaching $21.28$ on depth queries and $42.00$ overall
($42.07 \pm 0.61$ over three seeds against $36.07 \pm 1.75$), and reduces the seed-to-seed
variance of spatial queries by a factor of $4.6$.

The two results share a form. In the first, the specification failed to say \emph{what} a
program must return. In the second, it failed to say \emph{when} a capability applies.
Neither is a limitation of the model, and neither is visible until the model that had been
supplying the missing constraint is replaced.

\section{Future work}

\textbf{Complete the substitution register.} Building X-VLM (D9) would close the deviation
most likely to be costing accuracy on RefCOCO+, and would allow the attribute path to be
attributed rather than assumed.

\textbf{Resolve the candidate-set tension.} The detector threshold is tuned for detection
and is wrong for selection. A principled treatment --- a query-conditional candidate
budget, or a selection-aware detector --- would address a constraint that currently limits
both depth ordering and attribute checking, and would do so without the dataset artefact
that disqualified the naive cap.

\textbf{Extend the seed protocol} to RefCOCO+ and to the depth subset specifically, where
the present figures are single-run and small.

\textbf{Metric depth.} The primitives claim ordering only, because the datasets carry no
camera intrinsics. A dataset with intrinsics, or a metric depth model, would allow
\texttt{position\_3d} and \texttt{physical\_size} --- implemented but unused --- to be
evaluated, and would extend the API to relations that ordering cannot express.

\textbf{Generalise the specification finding.} The result here concerns one prompt and one
task. Whether other program-synthesis systems carry comparable unstated contracts is
testable by the same method: substitute the generator and measure what breaks.

\bibliographystyle{plain}
\bibliography{references}

\end{document}
